\documentclass{homework}
\author{Tashfeen, Ahmad}
\class{CSCI 3203: Tashfeen's Machine Learning I}
\title{Homework 6}
\address{
  Computer Science,
  Petree College of Arts \& Sciences,
  Oklahoma City University
}

\newcommand\synthetic{%
  \href{%
    https://tashfeen.org/raw/share/ml/synthetic.py%
  }{Download}%
}
\newcommand\sharepoint{%
  \href{%
    https://myokcuedu-my.sharepoint.com/:u:/g/personal/tashfeen_okcu_edu/EeqU1CH36W5Bp0YRgiRXxxUBuhzdI_JPec6CBlxsMBCtgA%
  }{\texttt{t10k-images-idx3-byte}}%
}
\newcommand\webarchive{%
  \href{%
    https://web.archive.org/web/20170130040655/https://yann.lecun.com/exdb/mnist/%
  }{the MNIST database of handwritten digits}%
}
\newcommand\linkedfile{%
  \href{%
    https://tashfeen.org/raw/share/ml/mnist.py%
  }{linked file}%
}

\begin{document} \maketitle

Let $x_{1}, \ldots, x_n$ be a set of instances in $\R^p$ and
$c_1, \ldots, c_K$ be a set of $K$ cluster centroids. Note that $c_i$
are also in $\R^p$. Let $\mathbb{I}_{ij} \in \{0, 1\}$ be a variable
such that $\mathbb{I}_{ij} = 1$ if $x_i$ is assigned to cluster
centred at $c_j$ and $\mathbb{I}_{ij} = 0$ otherwise. In layman's
terms, $\mathbb{I}_{ij}$ indicates whether $x_i$ is in the
$j^\text{th}$ cluster. The $K$--Means Clustering algorithm minimises
the following objective function,
\[
  \sum_{i=1}^n \sum_{j=1}^K \mathbb{I}_{ij} [(x_i - c_j)^T(x_i - c_j)]
\]

This objective function is essentially a measure of how far are the
members of each cluster from their respective centroids or means. The
$K$--Means Clustering algorithm is given bellow,

\begin{enumerate}
  \item Initialise centroids or means $c_j$ for $1 \leq j \leq K$
        randomly in $\R^p$ \ie $c_j \leftarrow x_j$.
  \item Assign all $x_i$ to the closest centroid $c_u$. I. e., for
        some $u$ and all $j$,
        \[
          \mathbf{if}\quad
          [(x_i - c_u)^T(x_i - c_u)]<[(x_i - c_j)^T(x_i - c_j)]
          \quad \mathbf{then}\quad
          \mathbb{I}_{iu}\leftarrow 1
        \]
  \item Adjust the centroids or means according to the new
        assignment in the last step,
        \[
          c_j \leftarrow \sum_{i=1}^{n}\mathbb{I}_{ij}\paren{\frac{x_i}{n}}
        \]
        Note that this is just assigning $c_j$ to the new mean of the
        $j^\text{th}$ cluster.
  \item Repeat from (2) till $c_j$ stop changing or the change
        is less than some $\ep$.
\end{enumerate}
This algorithm is in section 6.12.1 of Mitchell or optionally
algorithm 12.2 on page 523 of Hatie et al.

\question Listing \ref{code} plots the clusters shown in figure
\ref{clusters}. Give your best guesses for the optimal centroids
$c_j\in \R^2$ and $K=4$.

\begin{figure}[htbp]
  \begin{tabular}{p{0.60\textwidth}p{0.35\textwidth}}
    \begin{minipage}{.60\textwidth}
      \lstinputlisting[
        language={Python},
        caption={Synthetic 2D clusters. \synthetic.},
        label=code]
      {./code/synthetic.py}
    \end{minipage}
     &
    \begin{minipage}{.35\textwidth}
      \centering
      \includegraphics[width=\linewidth]{./media/clusters}
      \caption{Clusters}
      \label{clusters}
    \end{minipage}
  \end{tabular}
\end{figure}

\question Implement the $K$--Means Clustering algorithm for the data
in listing \ref{code}. Run the algorithm for $K \in \{1, 2, 3, 4\}$
and $\ep = 10^{-5}$. Give your final converged centroids plotted over
the clustered data. Plot each cluster with a distinct
colour\footnote{use different shades of grey if you have difficulty
  seeing colour}. At the end you should have four converged (with
different coloured clusters) plots for each value of $K$.


\question\label{lnk} Download the file \sharepoint. This is the binary
file containing 10,000 images of the test set from the
\webarchive. You can refer to their website for more information about
the data and its organisation. However, it is read and plotted for
your inspection and further usage to implement the $K$--Means
Clustering algorithm in the \linkedfile.

\begin{enumerate}
  \item Why do you think $K=10$ maybe a good choice for this data?
  \item What are the values of $n$ and $p$?
  \item What is represented by each of the $p$ features these data points?
  \item For $\ep = 3.5$, finish implementing the $K$--Means Clustering
        algorithm after the line indicated by the comment in the Python file
        given in question \ref{lnk}. The file has code that (when completed)
        should plot your 10 approximately converged cluster centroids as 28
        by 28 pixel images comprising of 2 rows and 5 columns. Give and
        discuss this plot. Why do you think the centroids $c_j$ look the way
        they do?
\end{enumerate}

\section*{Submission Instructions}
\begin{enumerate}
  \item Submit a PDF that answers the questions and contains all the
        plots that the assignment asks for.
  \item Submit your \texttt{kmeans.py}. This is your $K$--Means
        implementation for the code in listing \ref{code}.
  \item Submit your \texttt{mnsit.py}. This is your $K$--Means
        implementation for the MNSIT test set in question \ref{lnk}.
\end{enumerate}
\end{document}
